\section{总结}
本研究围绕保障性租赁住房的发展困境与对策展开系统分析，明确了保障性租赁住房的基本内涵及发展背景，
指出其作为解决新市民、青年人等群体住房问题的重要手段，具有政策性、非营利性和低租金等特征。
通过梳理我国保障性租赁住房的发展历程与政策演进，发现近年来在中央顶层设计推动下，相关制度框架逐步建立，取得了一定成效。

论文从资金、政策机制和居住品质三个维度总结了当前发展面临的主要困境，这些问题在一定程度上制约了保障性租赁住房的高质量发展。
并通过对北京、上海等典型城市的案例研究，总结经验，为全国推广提供了有益参考。
在对策建议层面，论文从健全政策制度、优化土地与资金供给、提升数字化治理水平、完善配套设施与社区融合等方面提出多项可行建议，
希望能推动我国保障性租赁住房实现可持续、高质量发展。